% ===== 字体（请使用 XeLaTeX 编译） =====
\usepackage{fontspec}
\defaultfontfeatures{Ligatures=TeX}

% ===== 页面与排版 =====
\usepackage[letterpaper,margin=1in,headheight=15pt]{geometry}
\usepackage{setspace}
\onehalfspacing
\usepackage{microtype}
\usepackage{indentfirst}
\setlength{\parindent}{2em}
\setlength{\parskip}{0.25em}

% ===== 数学、单位、表格与图片 =====
\usepackage{amsmath,amssymb,mathtools}
\usepackage{graphicx}
\graphicspath{{figures/}}
\usepackage{booktabs}
\usepackage{array}
\usepackage{caption}
\usepackage{subcaption}
\usepackage{enumitem}

% ===== 颜色 =====
\usepackage{xcolor}
\definecolor{UMblue}{HTML}{00274C}
\definecolor{UMmaize}{HTML}{FFCB05}

% ===== 文献 =====
\usepackage[round,authoryear]{natbib}

% ===== PDF 链接、书签与导航 =====
% hyperref 应尽量靠后加载；cleveref 必须在 hyperref 之后。
\usepackage{hyperref}
\usepackage{bookmark}
\usepackage[nameinlink,noabbrev]{cleveref}
\hypersetup{
  unicode=true,
  colorlinks=true,
  linkcolor=UMblue,
  citecolor=UMblue,
  urlcolor=UMblue,
  linktoc=all,
  bookmarksopen=true,
  bookmarksnumbered=true,
  pdfstartview=FitH
}

% ===== 页眉页脚 =====
\usepackage{fancyhdr}
\usepackage{lastpage}

\newcommand{\ContentsLink}{%
  \hyperlink{toc}{\textcolor{UMblue}{\bfseries Back to Contents}}%
}

\newcommand{\PageNavigation}{%
  \Acrobatmenu{PrevPage}{\textcolor{UMblue}{Previous}}%
  \quad\textcolor{gray}{|}\quad
  \hyperlink{toc}{\textcolor{UMblue}{Contents}}%
  \quad\textcolor{gray}{|}\quad
  \textcolor{black}{\thepage/\pageref*{LastPage}}%
  \quad\textcolor{gray}{|}\quad
  \Acrobatmenu{NextPage}{\textcolor{UMblue}{Next}}%
}

\fancypagestyle{reportstyle}{
  \fancyhf{}
  \fancyhead[L]{\small\nouppercase{\leftmark}}
  \fancyhead[R]{\small\ContentsLink}
  \fancyfoot[C]{\small\PageNavigation}
  \renewcommand{\headrulewidth}{0.4pt}
  \renewcommand{\footrulewidth}{0.4pt}
}

% chapter 首页默认使用 plain；重定义后仍保留全部导航。
\fancypagestyle{plain}{
  \fancyhf{}
  \fancyhead[L]{\small\nouppercase{\leftmark}}
  \fancyhead[R]{\small\ContentsLink}
  \fancyfoot[C]{\small\PageNavigation}
  \renewcommand{\headrulewidth}{0.4pt}
  \renewcommand{\footrulewidth}{0.4pt}
}

% ===== 目录深度与编号深度 =====
\setcounter{tocdepth}{2}    % 目录显示到 subsection
\setcounter{secnumdepth}{3} % 正文编号到 subsubsection

% ===== 常用命令示例 =====
\newcommand{\vect}[1]{\boldsymbol{#1}}
\newcommand{\dd}{\mathop{}\!\mathrm{d}}
